\addcontentsline{toc}{chapter}{Introduction}
\chapter*{Introduction}

These notes are intended as a guided meditation on the concept of the derivative. Here, ``derivative'' is interpreted broadly to refer to single variable derivatives, multivariable derivatives, and pushforwards of tangent vectors on manifolds. There are many exercises, because it seems to me the only way to understand concepts is to play with them. 


\section*{Prerequisites}

This text assumes basic familiarity with point-set topology (specifically, metric spaces, and topological properties of $\R$) and with linear algebra. At a few points, there may be some ideas from point-set topology and linear algebra that you have not encountered before. A sort of ``bare minimum'' exposition of some of these ideas is included in \cref{prelims}. 

\section*{How to use these notes}

If you're working through these notes and find one of the exercises challenging, that's great news! That means you'll have learned something when you solve it. But don't let the challenging exercises bog you down. Keep reading onwards, but do eventually come back to the exercises that gave you trouble. 

You are advised to \emph{not} spend time reading \cref{prelims} thoroughly before jumping into the main part of the text. When ideas from \cref{prelims} are invoked in the main part, a reference to the relevant part of \cref{prelims} is included. My suggestion is to only look at \cref{prelims} when you run into a reference to it, chasing references back as needed. 

Some sections are starred (\starred). This is intended to indicate that the section is more challenging, and perhaps not quite as important for the overall direction that these notes take. This is not to say they are not important; some very important facts are discussed in starred sections, and in fact unstarred sections do sometimes reference results in starred sections. But, if you find yourself struggling and need help deciding what's most important to focus on, focus on the unstarred sections. 

\section*{Suggestions for improvement}

This is a first draft, and there are bound to be many errors. Please be on the lookout for them! If you think you've found one, please share it with me. I'd also be very interested to hear about non-errors (sentences that are maybe correct but phrased confusingly, places where adding a picture could be helpful, etc).
